\documentclass[10pt,a4paper]{amsart}
\usepackage[margin=19mm]{geometry}
\usepackage{amsmath,amssymb,mathtools}
\usepackage[scr=rsfs,cal=euler]{mathalfa}
\usepackage{xfrac}
\usepackage[unicode,hidelinks,bookmarks=false]{hyperref}
\newcommand{\ie}{i.e.\ }
\newcommand{\cf}{cf.\ }
\newcommand{\resp}{respectively\ }
\newcommand{\Subsection}{\subsection}
\providecommand{\nobreakdash}{\nobreak\hbox{-}}
\setlength{\emergencystretch}{2em}
\multlinegap0pt
\usepackage{bm}
\usepackage{bbm}
\usepackage{dsfont}
\usepackage{xspace}
\usepackage{cancel}
\usepackage{mathtools}
\usepackage{amsthm}
\makeatletter
\@ifundefined{theorem}{\newtheorem{theorem}{Theorem}}{}
\@ifundefined{corollary}{\newtheorem{corollary}[theorem]{Corollary}}{}
\makeatother
\newcommand{\fqm}{\mathbb{F}_{q^m}}
\newcommand{\fqn}{\mathbb{F}_{q^n}}
\newcommand{\fq}{{\mathbb F}_{q}}
\newcommand{\la}{\langle}
\newcommand{\ra}{\rangle}
\title[Clubs in projective spaces and three-weight rank-metric code]{Clubs in projective spaces and three-weight rank-metric codes}
\author{Jonathan Mannaert, Paolo Santonastaso, Ferdinando Zullo}
\date{Source: arXiv:2508.00502v1 (CC BY 4.0)}
\begin{document}
\maketitle

\noindent\emph{Source and licence.} Clubs in projective spaces and three-weight rank-metric codes. Version arXiv:2508.00502v1 (CC BY 4.0), distributed under the Creative Commons Attribution 4.0 International licence, as recorded in the arXiv licence metadata for this exact version and shown on the abstract page https://arxiv.org/abs/2508.00502v1. Full rights notice: licenses/arxiv-2508.00502-CC-BY-4.0.txt.\par\smallskip

\noindent\textit{Editorial scope.} Counted unit: the Corollary whose source label is \texttt{None} in the article (source0.tex lines 988--996, source label \texttt{None}), delivered together with the article's own complete original proof of it (source0.tex lines 997--1045, one \texttt{proof} environment of 49 lines). No mathematics is written, repaired or re-derived: statement and proof are the article's own text line for line. Required context retained uncounted: none. Every definition, display and label of those lines is retained unchanged. Omitted from this package and not redistributed: 1--987, 1046--1434. Nothing inside a retained excerpt is omitted or altered: every retained body is delivered exactly as the article's lines are written, so no omission receipt of any kind accompanies this package. Citations: the retained excerpts carry no citation command at all, so no bibliography entry of the article is transcribed and no reference is redistributed. Reference adaptation: none was needed and none was made; every reference inside the delivered unit resolves inside it, so no unresolved reference or citation is produced. Printed slips kept literal: no printed word, symbol, hypothesis or number of the counted statement or of its proof was altered. Layout and class: the article's own class is \texttt{amsart} with packages including geometry, enumerate, amsmath, amssymb, latexsym, amsthm, color, xcolor, cancel, graphicx, cite, changes, etoolbox, hyperref; the wrapper is the lane's pinned \texttt{amsart} editorial wrapper at 10pt on A4, which restates the theorem-like environments the retained text needs and adds the article's own macro definitions that the retained text needs (\texttt{\textbackslash fq}, \texttt{\textbackslash fqm}, \texttt{\textbackslash fqn}, \texttt{\textbackslash la}, \texttt{\textbackslash ra}), with extra wrapper packages (bm, bbm, dsfont, xspace, cancel, mathtools). The delivered numbering is the unit's own. Assets: no retained span contains a figure environment, a graphics-inclusion command, a PDF-page inclusion, a table or a TikZ picture; no image file is copied into this package, the package declares no asset dependency and no image file is redistributed with it. Licence: version arXiv:2508.00502v1 of this article is distributed under the Creative Commons Attribution 4.0 International licence, as recorded in the arXiv licence metadata for this exact version and shown on the abstract page https://arxiv.org/abs/2508.00502v1. A full rights notice is delivered at licenses/arxiv-2508.00502-CC-BY-4.0.txt. Characters: every retained excerpt is pure ASCII, so the delivered engine, XeLaTeX with its default Latin Modern text font, reports no missing character.
\begin{corollary}
    Let $k$ and $m$ be two positive integers such that $k=2s+1$, for some positive integer $s$.
    Consider
    $$U=\{(\eta,x_1,x_1^{q}, x_2, x_2^{q}, \ldots, x_s,x_s^{q} )\mid x_1,\ldots,x_s\in \mathbb{F}_{q^m}, \eta\in S\},$$
    and
    $$\overline{U}=\{(x_1+\zeta,x_1^q,x_1^{q^2}, x_2, x_2^{q}, \ldots, x_s,x_s^{q} )\mid x_1,\ldots,x_s\in \mathbb{F}_{q^m}, \zeta\in T\},$$
    where $S$ and $T$ are $\fq$-subspaces of $\fqm$ with dimension $m-2$ such that there exist no $\alpha \in \fqn$ and $\sigma \in \mathrm{Aut}(\fqm)$ such that $S=\alpha T^\sigma$.
    We have that $U$ and $\overline{U}$ are $\mathrm{\Gamma L}(k,q^m)$-inequivalent.
\end{corollary}

\begin{proof}
    By contradiction, assume that $U$ and $\overline{U}$ are $\mathrm{\Gamma L}(k,q^m)$-equivalent. Therefore, there exist a matrix $A \in \mathrm{GL}(k,q^m)$ and $\sigma \in \mathrm{Aut}(\fqm)$ such that for every $u \in \overline{U}$ there exists $v \in U$ such that
    \[ A (u^\sigma)^T=v. \]
    Denote by $A_{i,j}$ the $(i,j)$-th entry of $A$. 
    We have that for every $x_1,\ldots,x_s\in \mathbb{F}_{q^m}, \zeta\in T$ there exist $y_1,\ldots,y_s\in \mathbb{F}_{q^m}, \eta\in S$ such that
    \[
    \left\{
    \begin{array}{llll}
    \eta =A_{1,1}(x_1^\sigma+\zeta^\sigma)+A_{1,2}(x_1^\sigma)^q+\ldots+A_{1,k}(x_s^\sigma)^q,\\
    y_1=A_{2,1}(x_1^\sigma+\zeta^\sigma)+A_{2,2}(x_1^\sigma)^q+\ldots+A_{2,k}(x_s^\sigma)^q,\\
    y_1^q=A_{3,1}(x_1^\sigma+\zeta^\sigma)+A_{3,2}(x_1^\sigma)^q+\ldots+A_{3,k}(x_s^\sigma)^q,\\
    \vdots\\
    y_s=A_{2s,1}(x_1^\sigma+\zeta^\sigma)+A_{2s,2}(x_1^\sigma)^q+\ldots+A_{2s,k}(x_s^\sigma)^q,\\
    y_1^q=A_{2s+1,1}(x_1^\sigma+\zeta^\sigma)+A_{2s+1,2}(x_1^\sigma)^q+\ldots+A_{2s+1,k}(x_s^\sigma)^q.
    \end{array}
    \right.
    \]
    In particular, since $\la (1,0,\ldots,0)\ra_{\fqm}$ is the only point of weight greater than one for both the linear sets $L_U$ and $L_{\overline{U}}$, it has to be fixed by the action of $A$ and $\sigma$. Hence, we have that $A_{1,1}\neq 0$ and for every $\zeta\in T$ there exists $\eta\in S$ such that 
    \[ A_{1,1} \zeta^\sigma=\eta,\]
    which implies that $A_{1,1}T^\sigma=S$, a contradiction.
    %for any $i \in [s]$, we have that
    %\[
    %\left\{
    %\begin{array}{llll}
    %y_i=A_{2i,1}(x_1+\zeta)+A_{2i,2}x_1^q+\ldots+A_{2i,k}x_s^q,\\
    %y_i^q=A_{2i+1,1}(x_1+\zeta)+A_{2i+1,2}x_1^q+\ldots+A_{2i+1,k}x_s^q,
    %\end{array}
    %\right.
    %\]
    %and by replacing the first equation into the second one we get
    %\[ A_{2i,1}^q(x_1^q+\zeta^q)+A_{2i,2}^qx_1^{q^2}+\ldots+A_{2i,k}^qx_s^{q^2}=A_{2i+1,1}(x_1+\zeta)+A_{2i+1,2}x_1^q+\ldots+A_{2i+1,k}x_s^q. \] 
    %The above equation must be valid for every $x_1,\ldots,x_s\in \mathbb{F}_{q^m}, \zeta\in T$ and so, looking at the above expression as a polynomial in $x_1,\ldots,x_s$, we get
    %\[
    %\left\{
    %\begin{array}{llll}
    %A_{2i,3}=0,\\
    %A_{2i+1,1}=0,\\
    %A_{2i,1}=0,\\
    %A_{2i+1,2}=0,\\
    %A_{2i+1,3}=A_{2i,2}^q,\\
    %A_{2i+1,4}=0,\\
    %A_{2i,5}=0,\\
    %A_{2i+1,5}=A_{2i,4}^q,\\
    %\vdots
    %\end{array}
    %\right.
    %\]
    %which implies in particular that
\end{proof}

\end{document}
